% !TEX program = xelatex
% Отчёт о лабораторной работе № 1 по дисциплине «Дискретная математика»
% «Генерация бинарного кода Грея. Мультимножества и операции над ними»
% Компиляция: xelatex Lab1_Multisets.tex (дважды — для содержания и ссылок)
% Рисунки fig_*.png строит make_figures.py из сценариев scenarios/*.txt; исходный код — в ../src
\documentclass[14pt,a4paper]{extarticle}

% ---------- Шрифты и язык ----------
\usepackage{fontspec}
\IfFontExistsTF{Times New Roman}
  {\setmainfont{Times New Roman}}
  {\setmainfont{Liberation Serif}}
\IfFontExistsTF{Courier New}
  {\setmonofont{Courier New}}
  {\setmonofont{Liberation Mono}}
\usepackage[english]{babel}
\babelprovide[import, main]{russian}
\usepackage{unicode-math}
\setmathfont{latinmodern-math.otf}
\IfFontExistsTF{Times New Roman}
  {\setmathfont[range={up/{latin,Latin,num}}]{Times New Roman}
   \setmathfont[range={it/{latin,Latin}}]{Times New Roman Italic}}
  {\setmathfont[range={up/{latin,Latin,num}}]{Liberation Serif}
   \setmathfont[range={it/{latin,Latin}}]{Liberation Serif Italic}}

% ---------- Поля и интервалы (ГОСТ 7.32-2017, п. 6.1.1) ----------
\usepackage[left=30mm,right=15mm,top=20mm,bottom=20mm,footskip=10mm]{geometry}
\usepackage{setspace}
\onehalfspacing
\setlength{\parindent}{1.25cm}
\usepackage{indentfirst}
\setlength{\parskip}{0pt}
\sloppy

% ---------- Заголовки ----------
\usepackage{titlesec}
\titleformat{\section}{\normalfont\bfseries}{\thesection}{0.6em}{}
\titleformat{\subsection}{\normalfont\bfseries}{\thesubsection}{0.6em}{}
\titlespacing*{\section}{\parindent}{0pt}{12pt}
\titlespacing*{\subsection}{\parindent}{12pt}{6pt}
\renewcommand{\thesection}{\arabic{section}}
\renewcommand{\thesubsection}{\thesection.\arabic{subsection}}

% Каждый раздел основной части — с новой страницы
\newcommand{\sectionbreak}{\clearpage}

% Структурные элементы без номера: с новой страницы, по центру, прописными
\newcommand{\structsection}[1]{%
  \clearpage
  \phantomsection
  \addcontentsline{toc}{section}{#1}%
  \begin{center}\textbf{#1}\end{center}\vspace{6pt}}

% ---------- Содержание ----------
\makeatletter
\renewcommand{\tableofcontents}{%
  {\centering\textbf{СОДЕРЖАНИЕ}\par}\vspace{12pt}%
  \@starttoc{toc}}
\renewcommand*{\l@section}{\@dottedtocline{1}{0pt}{1.5em}}
\renewcommand*{\l@subsection}{\@dottedtocline{2}{1.5em}{2.3em}}
\renewcommand{\@pnumwidth}{1.8em}
\makeatother

% ---------- Подписи рисунков и таблиц (ГОСТ 7.32-2017, пп. 6.5, 6.6) ----------
\usepackage{caption}
\DeclareCaptionLabelSeparator{gost}{~---~}
\captionsetup{labelsep=gost, font={normalsize}, skip=6pt}
\captionsetup[figure]{name=Рисунок, justification=centering}
\captionsetup[table]{name=Таблица, justification=raggedright, singlelinecheck=false, position=top}

% ---------- Таблицы, графика, формулы ----------
\usepackage{graphicx}
\usepackage{float}
\usepackage{array,tabularx,makecell,longtable}
\renewcommand{\arraystretch}{1.15}
\newcolumntype{C}{>{\centering\arraybackslash}X}
\newcolumntype{L}{>{\raggedright\arraybackslash}X}
\usepackage{amsmath}
\usepackage{enumitem}
\setlist{nosep, leftmargin=*, labelindent=\parindent, itemindent=0pt}
\newlist{gostlist}{enumerate}{1}
% Буквенная нумерация перечислений (без букв ё, з, й, о, ч, ъ, ы, ь)
\newcommand{\gostletter}[1]{\ifcase#1\or а\or б\or в\or г\or д\or е\or ж\or и\or к\or л\or м\or н\or п\fi}
\newcommand{\gostlettercnt}[1]{\gostletter{\value{#1}}}
\AddEnumerateCounter{\gostlettercnt}{\gostletter}{м}
\setlist[gostlist]{label=\gostlettercnt*), leftmargin=*, labelindent=\parindent, nosep}
\newlist{dashlist}{itemize}{1}
\setlist[dashlist]{label=--, leftmargin=*, labelindent=\parindent, nosep}
\newlist{steps}{enumerate}{1}
\setlist[steps]{label=\arabic*., leftmargin=*, labelindent=\parindent, nosep}

% Код выводится fvextra: listings под XeLaTeX путает порядок символов в кириллических комментариях
\usepackage{fvextra}
\fvset{numbers=left, numbersep=4pt, fontsize=\footnotesize, frame=single, breaklines=true,
       breaksymbolleft={}, breakanywhere=true, tabsize=4}
\usepackage{newfloat}
\DeclareFloatingEnvironment[name=Листинг, listname={Листинги}]{listing}
\captionsetup[listing]{justification=centering}
% \codefile[опции fancyvrb]{подпись}{метка}{файл}
\newcommand{\codefile}[4][]{%
  \par\medskip\captionof{listing}{#2}\label{#3}%
  \VerbatimInput[#1]{#4}\medskip}
% Символы из кода, которых нет в моноширинном шрифте
\usepackage{newunicodechar}
\newunicodechar{∩}{\ensuremath{\cap}}
\newunicodechar{∪}{\ensuremath{\cup}}
\newunicodechar{∆}{\ensuremath{\triangle}}
\newunicodechar{⊆}{\ensuremath{\subseteq}}
\newunicodechar{№}{\textnumero}
% В latinmodern-math нет глифа U+29F5, используем обратную косую черту
\AtBeginDocument{\renewcommand{\setminus}{\mathbin{\backslash}}}
\newcommand{\code}[1]{\texttt{#1}}

\usepackage[hidelinks]{hyperref}

\pagestyle{plain}

% ---------- Данные для титульного листа ----------
\newcommand{\StudentGroup}{5130201/50302}

% Обозначения
\newcommand{\kA}{k_A(x)}
\newcommand{\kB}{k_B(x)}
\newcommand{\kU}{k_U(x)}
\newcommand{\Kmax}{K_{\max}}

% Рисунок консоли: ширина по тексту, высота не больше 3/4 страницы
\newcommand{\consolefig}[3][0.95]{%
  \begin{figure}[H]
  \centering
  \includegraphics[width=#1\textwidth,height=0.74\textheight,keepaspectratio]{#2}
  \caption{#3}
  \label{fig:#2}
  \end{figure}}

\begin{document}
% ======================= ТИТУЛЬНЫЙ ЛИСТ =======================
\begin{titlepage}
\begin{singlespace}
\begin{center}
Министерство науки и высшего образования Российской Федерации\\[2pt]
Санкт-Петербургский политехнический университет Петра Великого\\[2pt]
Институт компьютерных наук и кибербезопасности\\[2pt]
Высшая школа технологий искусственного интеллекта

\vspace{5cm}

\textbf{ОТЧЁТ}\\[6pt]
\textbf{о лабораторной работе № 1}\\[12pt]
\textbf{«Генерация бинарного кода Грея.}\\
\textbf{Мультимножества и операции над ними»}\\[12pt]
по дисциплине «Дискретная математика»
\end{center}

\vspace{4cm}

\noindent
\begin{tabular}{@{}p{7.2cm}p{3.6cm}l@{}}
Выполнил\\
студент гр. \StudentGroup & \underline{\hspace{3cm}} & Р.\,Ю.~Дужнов\\[22pt]
Проверил\\
преподаватель & \underline{\hspace{3cm}} & А.\,В.~Востров\\
\end{tabular}

\vspace{1.2cm}
\noindent\hfill «\underline{\hspace{0.9cm}}» \underline{\hspace{3.2cm}} 2026 г.

\vfill
\begin{center}
Санкт-Петербург\\
2026
\end{center}
\end{singlespace}
\end{titlepage}
\setcounter{page}{2}

% ======================= СОДЕРЖАНИЕ =======================
\begin{singlespace}
\tableofcontents
\end{singlespace}

% ======================= ВВЕДЕНИЕ =======================
\structsection{ВВЕДЕНИЕ}

Требуется реализовать программу, которая генерирует бинарный код Грея заданной пользователем разрядности $n$ и заполняет им универсум мультимножеств. На основе универсума формируются два мультимножества $A$ и $B$; для каждого пользователь выбирает способ заполнения (вручную или автоматически) и задаёт мощность. Программа выводит результаты операций: объединение, пересечение, разность (оба варианта), симметрическая разность, дополнение (оба варианта), арифметические сумма, разность, произведение и деление. Требуются пользовательское меню и защита от некорректного ввода.

В работе приняты следующие уточнения постановки.
\begin{gostlist}
\item Универсум $U$~--- мультимножество: каждый код Грея $x$ входит в него с кратностью $\kU \geqslant 1$. Кратности задаются тремя способами: случайно от~1 до~$K$, вручную для каждого элемента или по заданной мощности $|U|$ со случайным распределением кратностей.
\item $A$ и $B$~--- подмультимножества $U$: $\kA \leqslant \kU$ для всех $x$. Мощность $|A|$ вводится пользователем или выбирается случайно; в обоих случаях $0 \leqslant |A| \leqslant |U|$, поэтому допустимо любое подмультимножество $U$, включая $\varnothing$ и сам $U$.
\item Некоммутативные арифметические операции (разность и деление) вычисляются в обоих вариантах, поэтому результатов тринадцать.
\item При $n = 0$ универсум пуст: $U = \{\}$, и все операции дают $\{\}$ (обоснование~--- в подразделе~\ref{sec:universe}).
\end{gostlist}

Входные данные программы перечислены в таблице~\ref{tab:input}. Пустая строка (только клавиша Enter) там, где она допустима, означает либо случайный выбор, либо отмену; в остальных местах она считается ошибкой.

\begin{table}[H]
\caption{Входные данные программы}
\label{tab:input}
\centering\small
\begin{tabularx}{\textwidth}{|L|l|L|}
\hline
\multicolumn{1}{|c|}{Величина} & \multicolumn{1}{c|}{Допустимые значения} & \multicolumn{1}{c|}{Пустая строка} \\ \hline
Пункт меню & целое от 0 до 6 & ошибка \\ \hline
Разрядность кода Грея $n$ & целое от 0 до 19 & ошибка \\ \hline
Способ задания кратностей $U$ & целое от 0 до 3 (0~--- отмена) & ошибка \\ \hline
Наибольшая кратность $K$ (способ~1) & целое от 1 до 100 & ошибка \\ \hline
Кратность $\kU$ (способ~2) & целое от 1 до 100 & отмена создания $U$ \\ \hline
Мощность $|U|$ (способ~3) & целое от $2^n$ до $100 \cdot 2^n$ & случайная мощность \\ \hline
Способ заполнения $A$, $B$ & целое от 0 до 2 (0~--- отмена) & ошибка \\ \hline
Мощность $|A|$, $|B|$ & целое от 0 до $|U|$ & случайная мощность \\ \hline
Код элемента (ручное заполнение) & строка из $n$ символов 0 и 1 & отмена заполнения \\ \hline
Добавляемая кратность & целое от 1 до $\min\{\kU - \kA,\ |A| - \text{набрано}\}$ & ошибка \\ \hline
\end{tabularx}
\end{table}

Результат считается достигнутым, если для любых допустимых входных данных программа выводит $U$, $A$, $B$ и тринадцать результатов операций в точном соответствии с определениями, а на любой некорректный ввод отвечает сообщением и повторным запросом, не завершаясь аварийно.

% ======================= 1 =======================
\section{Математическое описание}

\subsection{Множества}

Множество~--- совокупность различных элементов~[1]. Если все рассматриваемые множества являются подмножествами одного множества $U$, его называют универсумом. Операции над $A, B \subseteq U$~[1]:
\begin{align}
A \cup B &= \{x \mid x \in A \text{ или } x \in B\}, \label{eq:set-union}\\
A \cap B &= \{x \mid x \in A \text{ и } x \in B\}, \label{eq:set-inter}\\
A \setminus B &= \{x \mid x \in A \text{ и } x \notin B\} = A \cap \overline{B}, \label{eq:set-diff}\\
A \mathbin{\triangle} B &= (A \setminus B) \cup (B \setminus A), \label{eq:set-sym}\\
\overline{A} &= U \setminus A. \label{eq:set-compl}
\end{align}

Дополнение~\eqref{eq:set-compl} определено только относительно универсума. Тождества $A \setminus B = A \cap \overline{B}$ и~\eqref{eq:set-sym} переносятся на мультимножества и используются в реализации.

\subsection{Мультимножества}

Мультимножество $\hat{X}$ над множеством $X = \{x_1, \ldots, x_s\}$~--- совокупность элементов $X$, в которую $x_i$ входит $a_i \geqslant 0$ раз~[1]. Множество $X$ называют носителем, число $a_i$~--- кратностью (показателем) элемента $x_i$. Запись:
\begin{equation}
\hat{X} = \{x_1^{a_1}, \ldots, x_s^{a_s}\},
\end{equation}
элементы с нулевой кратностью не записываются. Мультимножество $A$ над $X$ однозначно задаётся функцией кратности $k_A \colon X \to \{0, 1, 2, \ldots\}$, а его мощность~--- сумма кратностей:
\begin{equation}
|A| = \sum_{x \in X} \kA. \label{eq:card}
\end{equation}
Мультимножество, все кратности которого равны 0 или 1, совпадает с обычным множеством.

Универсум $U$ в данной работе тоже мультимножество над носителем $X$ с кратностями $\kU \geqslant 1$. Мультимножество $A$ называется подмультимножеством $U$ ($A \subseteq U$), если
\begin{equation}
\kA \leqslant \kU \quad \text{для всех } x \in X. \label{eq:sub}
\end{equation}
Мультимножества $A$ и $B$ рассматриваются только как подмультимножества $U$. Из~\eqref{eq:card} и~\eqref{eq:sub} следует $0 \leqslant |A| \leqslant |U|$, причём каждое значение из этого отрезка достигается: подмультимножество мощности $m$ получается, если взять любые $m$ из $|U|$ «единиц кратности» универсума.

\subsection{Операции над мультимножествами}

Все операции выполняются поэлементно: кратность $x$ в результате зависит только от $\kA$, $\kB$ и $\kU$~[1,~4]. Функции кратности результатов приведены в таблице~\ref{tab:ops}.

\begin{table}[H]
\caption{Операции над мультимножествами}
\label{tab:ops}
\centering\small
\begin{tabularx}{\textwidth}{|L|c|L|}
\hline
\multicolumn{1}{|c|}{Операция} & Обозначение & \multicolumn{1}{c|}{Кратность $x$ в результате} \\ \hline
Объединение & $A \cup B$ & $\max\{\kA, \kB\}$ \\ \hline
Пересечение & $A \cap B$ & $\min\{\kA, \kB\}$ \\ \hline
Дополнение & $\overline{A}$ & $\kU - \kA$ \\ \hline
Разность & $A \setminus B$ & $\min\{\kA,\ \kU - \kB\}$ \\ \hline
Симметрическая разность & $A \mathbin{\triangle} B$ & $\max\bigl\{\min\{\kA, \kU - \kB\},\ \min\{\kB, \kU - \kA\}\bigr\}$ \\ \hline
Арифметическая сумма & $A + B$ & $\min\{\kA + \kB,\ \kU\}$ \\ \hline
Арифметическая разность & $A - B$ & $\max\{\kA - \kB,\ 0\}$ \\ \hline
Арифметическое произведение & $A \ast B$ & $\min\{\kA \cdot \kB,\ \kU\}$ \\ \hline
Арифметическое деление & $A / B$ & $\lfloor \kA / \kB \rfloor$ при $\kB > 0$; 0 при $\kB = 0$ \\ \hline
\end{tabularx}
\end{table}

Формула разности получена из $A \setminus B = A \cap \overline{B}$, симметрической разности~--- из~\eqref{eq:set-sym}. Арифметическая разность не зависит от универсума и отличается от теоретико-множественной: например, при $\kU = 5$, $\kA = 3$, $\kB = 1$ имеем $k_{A \setminus B}(x) = \min\{3, 4\} = 3$, но $k_{A - B}(x) = 2$. Сумма и произведение ограничены сверху кратностью $\kU$, иначе результат мог бы не быть подмультимножеством $U$ и к нему нельзя было бы применить дополнение.

\textbf{Утверждение.} Если $A, B \subseteq U$, то результат каждой операции из таблицы~\ref{tab:ops} является подмультимножеством $U$.

Действительно: $\max\{\kA, \kB\} \leqslant \kU$ и $\min\{\kA, \kB\} \leqslant \kU$ в силу~\eqref{eq:sub}; $0 \leqslant \kU - \kA \leqslant \kU$; разность и симметрическая разность не больше $\kU - \kB$ или $\kU - \kA$; сумма и произведение ограничены явно; $\max\{\kA - \kB, 0\} \leqslant \kA$; $\lfloor \kA/\kB \rfloor \leqslant \kA$ при $\kB \geqslant 1$.

Разность, арифметическая разность и деление некоммутативны, дополнение одноместно, поэтому вычисляются оба варианта: $A \setminus B$ и $B \setminus A$, $\overline{A}$ и $\overline{B}$, $A - B$ и $B - A$, $A / B$ и $B / A$. Вместе с $A \cup B$, $A \cap B$, $A \mathbin{\triangle} B$, $A + B$ и $A \ast B$ это тринадцать результатов.

Для проверки результатов используются следствия определений:
\begin{equation}
|\overline{A}| = |U| - |A|, \qquad |A \cup B| + |A \cap B| = |A| + |B|. \label{eq:check}
\end{equation}
Второе равенство верно поэлементно: $\max\{a, b\} + \min\{a, b\} = a + b$.

\subsection{Универсум на коде Грея}
\label{sec:universe}

Носитель $X$~--- множество всех $2^n$ двоичных слов длины $n$, записанных в порядке кода Грея (подраздел~\ref{sec:gray}); элемент универсума с номером $i$~--- $i$-е слово кода. Кратность каждого элемента в универсуме лежит в отрезке $[1, \Kmax]$, $\Kmax = 100$; поэтому
\begin{equation}
2^n \leqslant |U| \leqslant \Kmax \cdot 2^n. \label{eq:ucard}
\end{equation}
Кратности задаются одним из трёх способов:
\begin{gostlist}
\item случайно: $\kU = 1 + \xi_x$, где $\xi_x$ независимы и равновероятны на $\{0, \ldots, K - 1\}$, $1 \leqslant K \leqslant \Kmax$; тогда $\kU$ равновероятна на $[1, K]$, а математическое ожидание $|U|$ равно $2^n (K + 1)/2$;
\item вручную: пользователь вводит $\kU \in [1, \Kmax]$ для каждого из $2^n$ элементов;
\item по мощности: пользователь задаёт $|U|$ из~\eqref{eq:ucard} (или она выбирается равновероятно из этого отрезка); каждый элемент получает кратность 1, а остальные $|U| - 2^n$ единиц распределяются случайно между элементами так, чтобы ни одна кратность не превысила $\Kmax$ (алгоритм подраздела~\ref{sec:random}, применённый к мультимножеству, в котором все кратности равны $\Kmax - 1$).
\end{gostlist}

При $K = 1$ (способ 1) или $|U| = 2^n$ (способ 3) все кратности равны 1, $U$ и его подмультимножества становятся обычными множествами, а арифметические сумма, разность и произведение совпадают с объединением, разностью и пересечением.

\textbf{Случай $n = 0$.} Двоичное слово длины 0 не содержит ни одного разряда. Такое слово нельзя ни показать в таблице универсума, ни ввести с клавиатуры как код элемента (пустая строка в программе означает отмену ввода). Поэтому в программе принято, что при $n = 0$ носитель пуст: $X = \varnothing$, $U = \{\}$, $|U| = 0$. Из~\eqref{eq:sub} следует, что единственное подмультимножество пустого универсума~--- пустое, и допустимая мощность $A$ и $B$ равна только 0. Все функции кратности из таблицы~\ref{tab:ops} определены на пустом носителе, поэтому каждая из тринадцати операций даёт $\{\}$ с мощностью 0; соотношения~\eqref{eq:check} выполняются: $0 = 0 - 0$ и $0 + 0 = 0 + 0$. Программа проходит этот случай теми же функциями, что и общий (сценарий~5 в подразделе~\ref{sec:s5}).

\subsection{Алгоритм генерации бинарного кода Грея}
\label{sec:gray}

Бинарный код Грея порядка $n$~--- последовательность всех $2^n$ двоичных слов длины $n$, в которой соседние слова, а также последнее и первое, различаются ровно в одном разряде~[1].

Слово рассматривается как битовая шкала $B = (B[n], \ldots, B[2], B[1])$; разряды нумеруются справа, $B[1]$~--- младший. На шаге $i$ инвертируется разряд с номером $Q(i)$, где
\begin{equation}
Q(i) = \text{число нулей в конце двоичной записи } i, \text{ увеличенное на } 1. \label{eq:q}
\end{equation}
$Q(i)$ вычисляется так: $q := 1$, $j := i$; пока $j$ чётно, $j := j/2$ и $q := q + 1$; результат~--- $q$.

\textbf{Вход:} целое $n \geqslant 1$.

\textbf{Выход:} последовательность из $2^n$ шкал длины $n$.

\begin{steps}
\item Положить $B[k] := 0$ для $k = 1, \ldots, n$ и выдать $B$.
\item Для $i = 1, 2, \ldots, 2^n - 1$: $p := Q(i)$; $B[p] := 1 - B[p]$; выдать $B$.
\end{steps}

При $i < 2^n$ двоичная запись $i$ оканчивается не более чем $n - 1$ нулями, поэтому $1 \leqslant Q(i) \leqslant n$ и номер разряда всегда допустим. Протокол работы при $n = 3$~--- в таблице~\ref{tab:gray}.

\begin{table}[H]
\caption{Протокол алгоритма при $n = 3$}
\label{tab:gray}
\centering
\begin{tabularx}{0.85\textwidth}{|C|C|C|C|}
\hline
Шаг $i$ & Двоичная запись $i$ & $p = Q(i)$ & Шкала $B$ \\ \hline
--- & --- & --- & 000 \\ \hline
1 & 001 & 1 & 001 \\ \hline
2 & 010 & 2 & 011 \\ \hline
3 & 011 & 1 & 010 \\ \hline
4 & 100 & 3 & 110 \\ \hline
5 & 101 & 1 & 111 \\ \hline
6 & 110 & 2 & 101 \\ \hline
7 & 111 & 1 & 100 \\ \hline
\end{tabularx}
\end{table}

\textbf{Корректность.} Соседние шкалы различаются одним разрядом по построению. Докажем, что все $2^n$ шкал различны и последняя отличается от первой одним разрядом. Обозначим $S_k = \bigl(Q(1), \ldots, Q(2^k - 1)\bigr)$~--- последовательность номеров инвертируемых разрядов; она не зависит от $n$. Для $0 < j < 2^k$ числа $2^k + j$ и $j$ оканчиваются одинаковым числом нулей, а $Q(2^k) = k + 1$; значит,
\begin{equation}
S_{k+1} = \bigl(S_k,\ k + 1,\ S_k\bigr).
\end{equation}
Индукцией по $k$ отсюда следует, что $S_k$ совпадает со своим обращением (палиндром): $S_1 = (1)$, и обращение $(S_k, k{+}1, S_k)$ равно $(S_k^R, k{+}1, S_k^R) = (S_k, k{+}1, S_k)$.

Пусть при длине $k$ алгоритм выдаёт шкалы $g_1, \ldots, g_m$, $m = 2^k$, и утверждение для них верно. При длине $k + 1$ первые $m - 1$ шагов используют $S_k$ и не трогают старший разряд $k + 1$: получаются $0g_1, \ldots, 0g_m$. Шаг $m$ инвертирует старший разряд: $1g_m$. Оставшиеся $m - 1$ шагов снова применяют $S_k$, то есть, так как $S_k$~--- палиндром, обращение $S_k$, и проходят шкалы в обратном порядке: $1g_{m-1}, \ldots, 1g_1$. Итог~--- $(0g_1, \ldots, 0g_m, 1g_m, \ldots, 1g_1)$, отражённый код Грея. Шкалы внутри каждой половины различны по предположению индукции, половины различаются старшим разрядом, последняя шкала $1g_1$ отличается от первой $0g_1$ только старшим разрядом. База $k = 1$: шкалы 0 и 1. Получаемая последовательность совпадает с замкнутой формой $g(i) = i \oplus \lfloor i/2 \rfloor$; программа проверялась на совпадение при $n = 1, \ldots, 16$.

\textbf{Трудоёмкость.} Алгоритм делает $2^n - 1$ шагов; на шаге $i$ вычисление $Q(i)$ требует $Q(i)$ проверок чётности. Так как
\[
\sum_{i=1}^{2^n - 1} Q(i) = (2^n - 1) + \sum_{j=1}^{n} \Bigl\lfloor \frac{2^n - 1}{2^j} \Bigr\rfloor = 2^{n+1} - n - 2,
\]
общее время~--- $O(2^n)$, то есть $O(1)$ в среднем на слово. Память~--- $O(2^n)$ машинных слов при хранении каждой шкалы целым числом.

\subsection{Случайный выбор подмультимножества заданной мощности}
\label{sec:random}

Автоматическое заполнение $A$ мощности $m$ и способ 3 задания универсума сводятся к одной задаче: дано мультимножество-ограничение $V$ (это $U$ или «свободные места» универсума); выбрать случайное $A \subseteq V$ с $|A| = m$, $0 \leqslant m \leqslant |V|$.

Используется урновая модель. Элемент $x$ представлен $k_V(x)$ одинаковыми шарами с меткой $x$, всего шаров $N = |V|$. Из урны без возвращения извлекаются $m$ шаров, и $\kA$~--- число извлечённых шаров с меткой $x$. Выбор $m$ шаров из $N$ выполняется алгоритмом выборки~S~[2]: шары просматриваются по очереди, и текущий шар берётся с вероятностью
\begin{equation}
P = \frac{m - s}{N - t}, \label{eq:algs}
\end{equation}
где $t$~--- число уже просмотренных шаров, $s$~--- число уже взятых.

\textbf{Вход:} $V$ с кратностями $k_V(x)$, $N = |V|$; целое $m$.

\textbf{Выход:} $A \subseteq V$, $|A| = m$.

\begin{steps}
\item Если $m > N$, завершить с ошибкой.
\item Положить $\kA := 0$ для всех $x$, $\textit{need} := m$, $\textit{left} := N$.
\item Для элементов $x$ по порядку, пока $\textit{need} > 0$: повторить $k_V(x)$ раз, пока $\textit{need} > 0$: если $\textit{need} = \textit{left}$ или случайное целое $r \in [0, \textit{left} - 1]$ меньше $\textit{need}$, то $\kA := \kA + 1$ и $\textit{need} := \textit{need} - 1$; затем $\textit{left} := \textit{left} - 1$.
\item Вернуть $A$.
\end{steps}

Условие $r < \textit{need}$ при равновероятном $r$ выполняется с вероятностью $\textit{need}/\textit{left}$, что совпадает с~\eqref{eq:algs}.

\textbf{Корректность.} Всегда $0 \leqslant \textit{need} \leqslant \textit{left}$: при $\textit{need} = \textit{left}$ шар берётся обязательно. Поэтому к концу просмотра взято ровно $m$ шаров, $|A| = m$, а из каждого элемента взято не больше $k_V(x)$ шаров, то есть $A \subseteq V$. Вероятность выбрать заданный набор из $m$ шаров равна произведению вероятностей~\eqref{eq:algs} по всем просмотренным шарам; числители выбранных шаров дают $m!$, числители невыбранных~--- $(N - m)!$, знаменатели~--- $N!$:
\begin{equation}
P(\text{набор}) = \frac{m!\,(N - m)!}{N!} = \binom{N}{m}^{-1},
\end{equation}
то есть все $\binom{N}{m}$ наборов шаров равновероятны. Следовательно, $\kA$ имеет гипергеометрическое распределение, и
\begin{equation}
\mathrm{M}\,\kA = m \cdot \frac{k_V(x)}{N}:
\end{equation}
элементы с большей кратностью в универсуме получают в среднем пропорционально большую кратность в $A$.

\textbf{Трудоёмкость.} Просматривается не больше $N = |V|$ шаров, на каждый~--- одно случайное число; после того как $\textit{need}$ обратится в 0 или сравняется с $\textit{left}$, случайные числа не нужны. Время~--- $O(|V| + 2^n)$, без холостых выборок и с гарантированной верхней границей. При $|V| \leqslant \Kmax \cdot 2^n$ это $O(\Kmax \cdot 2^n)$; именно эта величина определяет наибольшую допустимую разрядность (подраздел~\ref{sec:limits}).

% ======================= 2 =======================
\section{Особенности реализации}

\subsection{Средства разработки и структура программы}

Язык C++14, только стандартная библиотека; компилятор GCC 6.3.0 (MinGW, 32-разрядная сборка), ОС Windows~11. Сборка:
\begin{center}
\code{g++ -std=c++14 -O2 -Isrc src/main.cpp src/ui/*.cpp src/core/*.cpp -o lab1}
\end{center}
или через CMake (файл \code{CMakeLists.txt}, цели \code{lab1} и \code{bench}). Предупреждений при \code{-Wall -Wextra} нет.

Код разделён на два слоя: \code{core}~--- математическая часть, не работающая с консолью; \code{ui}~--- ввод и меню. Зависимости направлены только от \code{ui} к \code{core}. Файлы перечислены в таблице~\ref{tab:files}, текст программы~--- в приложении~А.

\begin{table}[H]
\caption{Файлы программы}
\label{tab:files}
\centering\small
\begin{tabularx}{\textwidth}{|l|L|}
\hline
\multicolumn{1}{|c|}{Файл} & \multicolumn{1}{c|}{Назначение} \\ \hline
\code{src/core/GrayCode.hpp, .cpp} & генерация кода Грея, класс \code{CodeSet} (носитель) \\ \hline
\code{src/core/Multiset.hpp, .cpp} & класс \code{Multiset}, операции, случайный выбор подмультимножества, вывод \\ \hline
\code{src/core/Universe.hpp, .cpp} & три способа построения универсума \\ \hline
\code{src/core/Random.hpp, .cpp} & генератор случайных чисел \\ \hline
\code{src/ui/Input.hpp, .cpp} & защищённый построчный ввод \\ \hline
\code{src/ui/Menu.hpp, .cpp} & класс \code{Menu}: состояние программы и меню \\ \hline
\code{src/main.cpp} & точка входа, ключ \code{--seed} \\ \hline
\code{tools/bench.cpp} & замер времени для выбора наибольшей разрядности \\ \hline
\end{tabularx}
\end{table}

Все величины задачи (разрядность, номера, кратности, мощности, коды) неотрицательны и хранятся в типе \code{size\_t}. В использованной 32-разрядной сборке \code{size\_t} занимает 4 байта (наибольшее значение $2^{32} - 1 = 4\,294\,967\,295$). Наибольшая мощность универсума $100 \cdot 2^{19} = 52\,428\,800$, наибольшее произведение кратностей $100 \cdot 100$, наибольшее вводимое число $999\,999\,999$~--- переполнение невозможно. Вычитание беззнаковых величин выполняется только тогда, когда уменьшаемое заведомо не меньше: в дополнении это гарантирует проверка $A \subseteq U$, в арифметической разности~--- сравнение $\kA > \kB$, в ручном заполнении~--- инвариант $\kA \leqslant \kU$.

\subsection{Ограничения разрядности и вывода}
\label{sec:limits}

\textbf{Наибольшая разрядность $n_{\max} = 19$} (константа \code{MAX\_DEPTH}). Критерий выбора: любое действие пользователя должно завершаться быстрее 1~с. Это граница времени отклика, при которой пользователь ещё воспринимает работу с программой как непрерывную~[3]; при большей задержке интерактивное меню «зависает» в восприятии пользователя.

Самое долгое действие~--- случайный выбор подмультимножества (подраздел~\ref{sec:random}): его время пропорционально $|U| \leqslant \Kmax \cdot 2^n$. Остальные действия линейны по $2^n$ с малой константой. Время измерено утилитой \code{tools/bench.cpp} в наихудшем случае: все $\kU = \Kmax = 100$, мощность $|A| = |U|/2$ (при ней алгоритм просматривает почти все $|U|$ шаров), лучший из трёх запусков. Машина: Intel Core i7-13700KF, 32~ГБ ОЗУ; сборка \code{-O2}. Результаты~--- в таблице~\ref{tab:bench}.

\begin{table}[H]
\caption{Время действий, мс (наихудший случай, $\Kmax = 100$)}
\label{tab:bench}
\centering\small
\begin{tabularx}{\textwidth}{|c|C|C|C|C|C|}
\hline
$n$ & Код Грея & $U$ случайно (способ 1) & $U$ по мощности (способ 3) & Заполнение $A$ & 13 операций \\ \hline
16 & 0,0 & 0,0 & 72 & 72 & 2,0 \\ \hline
17 & 0,0 & 0,6 & 144 & 145 & 3,6 \\ \hline
18 & 0,5 & 1,5 & 288 & 381 & 12 \\ \hline
\textbf{19} & 1,0 & 6,0 & \textbf{873} & \textbf{877} & 23 \\ \hline
20 & 3,0 & 13 & 1752 & 1746 & 49 \\ \hline
21 & 7,0 & 25 & 3551 & 3529 & 125 \\ \hline
\end{tabularx}
\end{table}

Время растёт примерно вдвое на каждый разряд, так как вдвое растёт $|U|$. При $n = 19$ наибольшее время 0,88~с укладывается в 1~с, при $n = 20$ оно уже 1,75~с. Поэтому $n_{\max} = 19$. Память не ограничивает: при $n = 19$ один объект \code{Multiset} занимает $2^{19} \cdot 4$~байт $= 2$~МиБ, носитель~--- 4~МиБ.

Независимо от критерия времени есть предел типа: $2^n$ должно помещаться в \code{size\_t}, поэтому \code{generate\_gray\_code} отвергает $n > 31$ (константа \code{MAX\_TYPE\_DEPTH}). \code{MAX\_DEPTH}~--- ограничение интерфейса, \code{MAX\_TYPE\_DEPTH}~--- ограничение ядра; утилита замера пользуется ядром и потому может измерять $n > 19$.

\textbf{Наибольшая кратность $\Kmax = 100$} (константа \code{MAX\_MULTIPLICITY}): трёхзначные кратности умещаются в столбцы таблиц. От $\Kmax$ зависит $n_{\max}$: время автозаполнения пропорционально $\Kmax \cdot 2^n$, поэтому при изменении $\Kmax$ замер нужно повторить.

\textbf{Ограничения вывода.} Таблица выводит не больше 32 строк (\code{TABLE\_LIMIT}; при $n \leqslant 5$~--- полностью), запись мультимножества~--- не больше 16 элементов (\code{PRINT\_LIMIT}; при $n \leqslant 4$~--- всегда полностью). Дальше выводится пометка, сколько строк или элементов не показано, а мощности выводятся всегда полностью. При $n = 19$ полный вывод одного результата занял бы около $2^{19}$ записей и был бы необозрим.

\subsection{Исключения}

Ошибки пользователя исключений не вызывают: программа выводит сообщение и повторяет запрос. Исключения сигнализируют о нарушении внутренних ограничений, которое возможно только при ошибке в самой программе, и о конце потока ввода (таблица~\ref{tab:exc}).

\begin{table}[H]
\caption{Классы исключений}
\label{tab:exc}
\centering\small
\begin{tabularx}{\textwidth}{|l|l|L|}
\hline
\multicolumn{1}{|c|}{Класс} & \multicolumn{1}{c|}{Родитель} & \multicolumn{1}{c|}{Условие} \\ \hline
\code{GrayCodeError} & \code{std::runtime\_error} & общий родитель ошибок кода Грея \\ \hline
\code{InvalidDepthError} & \code{GrayCodeError} & $n > $ \code{MAX\_TYPE\_DEPTH} \\ \hline
\code{CodeIndexError} & \code{GrayCodeError} & номер кода вне носителя \\ \hline
\code{MultisetError} & \code{std::runtime\_error} & общий родитель ошибок мультимножеств \\ \hline
\code{MultisetIndexError} & \code{MultisetError} & номер элемента вне носителя \\ \hline
\code{CarrierMismatchError} & \code{MultisetError} & операнды построены над разными носителями \\ \hline
\code{NotSubmultisetError} & \code{MultisetError} & операнд не является подмультимножеством $U$ \\ \hline
\code{CardinalityOverflowError} & \code{MultisetError} & требуемая мощность больше мощности ограничения \\ \hline
\code{InvalidMultiplicityError} & \code{MultisetError} & $K \notin [1, \Kmax]$ \\ \hline
\code{UniverseCardinalityError} & \code{MultisetError} & $|U|$ вне~\eqref{eq:ucard} \\ \hline
\code{InputClosedError} & \code{std::runtime\_error} & поток ввода закончился \\ \hline
\end{tabularx}
\end{table}

\subsection{Модуль случайных чисел}

Источник случайности~--- генератор \code{std::mt19937} (вихрь Мерсенна, 32-разрядные значения); равновероятность обеспечивает \code{std::uniform\_int\_distribution}~[5]. Стандартная \code{std::rand} не подходит: в использованной сборке $\text{RAND\_MAX} = 32\,767$, и при $|U| > 32\,768$ часть шаров в алгоритме подраздела~\ref{sec:random} была бы недостижима.

\begin{dashlist}
\item \code{unsigned make\_seed()}~--- возвращает зерно: \code{std::random\_device} \textit{исключающее ИЛИ} текущее время. В MinGW \code{std::random\_device} детерминирован (проверено: два запуска выдают одинаковые числа), поэтому без времени все запуски давали бы одни и те же универсумы.
\item \code{void seed\_random(unsigned seed)}~--- задаёт начальное состояние генератора.
\item \code{size\_t random\_below(size\_t bound)}~--- вход: $\text{bound} > 0$; выход: равновероятное целое из $[0, \text{bound} - 1]$.
\end{dashlist}

Функция \code{main} выводит использованное зерно; запуск \code{lab1 --seed N} повторяет все случайные числа. Так получены воспроизводимые сценарии раздела~3.

\subsection{Генерация кода Грея и носитель}

Шкала $B$ хранится одним целым числом: разряд $B[p]$~--- бит числа с номером $p - 1$ (нумерация справа с нуля), инвертирование~--- \code{code \^{}= 1 << (p - 1)}. Реализация~--- в листинге~\ref{lst:gray}.

\codefile[firstline=4, lastline=35, firstnumber=4]{Генерация кода Грея (\code{GrayCode.cpp})}{lst:gray}{../src/core/GrayCode.cpp}

\begin{dashlist}
\item \code{static size\_t changed\_bit(size\_t i)}~--- функция $Q$~\eqref{eq:q}; вход: $i > 0$; выход: $Q(i)$. Видна только в \code{GrayCode.cpp}.
\item \code{std::vector<size\_t> generate\_gray\_code(size\_t depth)}~--- вход: $n$; выход: $2^n$ кодов в порядке построения; при $n = 0$~--- пустая последовательность; при $n > 31$~--- исключение \code{InvalidDepthError}.
\item \code{std::string code\_to\_string(size\_t code, size\_t depth)}~--- выход: строка из $n$ символов 0 и 1, старший разряд слева.
\end{dashlist}

Класс \code{CodeSet}~--- носитель мультимножеств (таблицы~\ref{tab:codeset-f}, \ref{tab:codeset-m}). Поле \code{positions}~--- обратная таблица: $\code{positions}[c]$~--- номер кода $c$ в последовательности. Так как код Грея содержит все $2^n$ слов, номер есть у каждого слова длины $n$, и поиск элемента по введённой строке выполняется за $O(n)$ (перевод строки в число и одно обращение к массиву) вместо $O(n \cdot 2^n)$ при просмотре всех кодов.

\begin{table}[H]
\caption{Поля класса \code{CodeSet}}
\label{tab:codeset-f}
\centering\small
\begin{tabularx}{\textwidth}{|l|l|L|}
\hline
\multicolumn{1}{|c|}{Поле} & \multicolumn{1}{c|}{Тип} & \multicolumn{1}{c|}{Назначение} \\ \hline
\code{depth} & \code{size\_t} & разрядность $n$ \\ \hline
\code{codes} & \code{std::vector<size\_t>} & \code{codes[i]}~--- $i$-й код Грея \\ \hline
\code{positions} & \code{std::vector<size\_t>} & \code{positions[c]}~--- номер кода $c$ \\ \hline
\end{tabularx}
\end{table}

\begin{table}[H]
\caption{Методы класса \code{CodeSet}}
\label{tab:codeset-m}
\centering\small
\begin{tabularx}{\textwidth}{|l|L|L|}
\hline
\multicolumn{1}{|c|}{Метод} & \multicolumn{1}{c|}{Вход} & \multicolumn{1}{c|}{Выход} \\ \hline
\code{CodeSet(depth)} & разрядность $n$ & носитель из $2^n$ кодов (пустой при $n = 0$) \\ \hline
\code{get\_depth()} & --- & $n$ \\ \hline
\code{size()} & --- & число элементов $2^n$ (0 при $n = 0$) \\ \hline
\code{get\_code(i)} & номер $i$ & строка кода; \code{CodeIndexError} при $i \geqslant 2^n$ \\ \hline
\code{find(code, index)} & строка; переменная для номера & \code{true} и номер в \code{index}, если строка~--- слово длины $n$ из 0 и 1; иначе \code{false} \\ \hline
\end{tabularx}
\end{table}

\subsection{Класс Multiset}

Класс \code{Multiset}~--- мультимножество над носителем, представленное функцией кратности. Объектами этого класса являются и универсум $U$, и мультимножества $A$, $B$, и все результаты операций: универсум отличается только тем, что его кратности не меньше 1 и он передаётся операциям, которым нужна $\kU$. Поля~--- в таблице~\ref{tab:ms-f}.

\begin{table}[H]
\caption{Поля класса \code{Multiset}}
\label{tab:ms-f}
\centering\small
\begin{tabularx}{\textwidth}{|l|l|L|}
\hline
\multicolumn{1}{|c|}{Поле} & \multicolumn{1}{c|}{Тип} & \multicolumn{1}{c|}{Назначение} \\ \hline
\code{carrier} & \code{std::shared\_ptr<const CodeSet>} & общий носитель; объект носителя живёт, пока на него ссылается хотя бы одно мультимножество \\ \hline
\code{counts} & \code{std::vector<size\_t>} & \code{counts[i]}~--- кратность $i$-го кода Грея \\ \hline
\end{tabularx}
\end{table}

Кратности хранятся по номерам элементов носителя, поэтому поиск элементов в операциях не нужен, и каждая операция выполняется за $O(2^n)$ одним проходом (разность~--- двумя, симметрическая разность~--- пятью, так как они собраны из других операций). Совместное владение носителем через \code{shared\_ptr} исключает висячие указатели: при пересоздании универсума старые мультимножества продолжают ссылаться на свой носитель. Принадлежность одному носителю проверяется сравнением указателей.

Методы доступа~--- в таблице~\ref{tab:ms-m}. Закрытый метод \code{check\_same\_carrier} вызывается каждой двуместной операцией, \code{check\_within}~--- каждой операцией, которой передаётся универсум.

\begin{table}[H]
\caption{Методы доступа класса \code{Multiset}}
\label{tab:ms-m}
\centering\small
\begin{tabularx}{\textwidth}{|l|p{3.2cm}|L|}
\hline
\multicolumn{1}{|c|}{Метод} & \multicolumn{1}{c|}{Вход} & \multicolumn{1}{c|}{Выход} \\ \hline
\code{Multiset(carrier)} & носитель & мультимножество с нулевыми кратностями \\ \hline
\code{size()} & --- & число элементов носителя \\ \hline
\code{get\_count(i)} & номер & кратность; \code{MultisetIndexError} \\ \hline
\code{set\_count(i, count)} & номер, кратность & кратность изменена; \code{MultisetIndexError} \\ \hline
\code{cardinality()} & --- & мощность по~\eqref{eq:card} \\ \hline
\code{is\_submultiset\_of(other)} & мультимножество & истина, если выполнено~\eqref{eq:sub} \\ \hline
\code{check\_same\_carrier(other)} & мультимножество & \code{CarrierMismatchError}, если носители разные \\ \hline
\code{check\_within(universe)} & универсум & то же, плюс \code{NotSubmultisetError}, если объект не подмультимножество $U$ \\ \hline
\end{tabularx}
\end{table}

Методы операций (таблица~\ref{tab:ms-ops}): первый операнд $A$~--- сам объект, второй $B$~--- параметр; каждый метод возвращает новый объект и не меняет операнды. Метод объединения назван \code{unite}, так как \code{union}~--- ключевое слово C++.

\begin{table}[H]
\caption{Методы операций класса \code{Multiset}}
\label{tab:ms-ops}
\centering\small
\begin{tabularx}{\textwidth}{|l|c|L|}
\hline
\multicolumn{1}{|c|}{Метод} & Результат & \multicolumn{1}{c|}{Реализация} \\ \hline
\code{unite(b)} & $A \cup B$ & \code{max} по элементам \\ \hline
\code{intersect(b)} & $A \cap B$ & \code{min} по элементам \\ \hline
\code{complement(u)} & $\overline{A}$ & $\kU - \kA$ после проверки $A \subseteq U$ \\ \hline
\code{difference(b, u)} & $A \setminus B$ & $A \cap \overline{B}$ методами \code{intersect} и \code{complement} \\ \hline
\code{symmetric\_difference(b, u)} & $A \mathbin{\triangle} B$ & $(A \setminus B) \cup (B \setminus A)$ методами \code{difference} и \code{unite} \\ \hline
\code{arithmetic\_sum(b, u)} & $A + B$ & $\min\{\kA + \kB, \kU\}$ \\ \hline
\code{arithmetic\_difference(b)} & $A - B$ & $\kA - \kB$ при $\kA > \kB$, иначе 0 \\ \hline
\code{arithmetic\_product(b, u)} & $A \ast B$ & $\min\{\kA \cdot \kB, \kU\}$ \\ \hline
\code{arithmetic\_division(b)} & $A / B$ & $\lfloor \kA / \kB \rfloor$ при $\kB \ne 0$, иначе 0 \\ \hline
\end{tabularx}
\end{table}

Функция \code{random\_submultiset(bound, cardinality)} реализует алгоритм подраздела~\ref{sec:random} (листинг~\ref{lst:random}); вход: ограничение $V$ и мощность $m$; выход: случайное $A \subseteq V$, $|A| = m$; при $m > |V|$~--- \code{CardinalityOverflowError}. Она объявлена другом класса, чтобы работать с \code{counts} напрямую.

\codefile[firstline=135, lastline=156, firstnumber=135]{Случайный выбор подмультимножества (\code{Multiset.cpp})}{lst:random}{../src/core/Multiset.cpp}

Вывод: \code{print\_multiset(out, multiset, limit)} записывает в поток элементы с ненулевой кратностью в виде \code{\{000(2), 011(1)\}}, не больше \code{limit} элементов, и пометку «\code{... ещё N эл.}»; пустое мультимножество записывается как \code{\{\}}. Оператор \code{<<} выводит мультимножество полностью.

Деструкторы, конструкторы копирования и перемещения не объявлены явно: классы содержат только \code{std::vector}, \code{std::shared\_ptr} и числа, и сгенерированные компилятором функции корректны.

\subsection{Построение универсума}

Модуль \code{Universe} реализует способы 1 и 3 подраздела~\ref{sec:universe}; способ 2 (ручной) выполняет меню, записывая введённые числа методом \code{set\_count}.
\begin{dashlist}
\item \code{min\_universe\_cardinality(size)} и \code{max\_universe\_cardinality(size)}~---\\ границы~\eqref{eq:ucard}: $2^n$ и $\Kmax \cdot 2^n$.
\item \code{random\_universe(carrier, K)}~--- вход: носитель, $K \in [1, \Kmax]$; выход: универсум с $\kU = 1 + \code{random\_below}(K)$; иначе \code{InvalidMultiplicityError}.
\item \code{universe\_with\_cardinality(carrier, m)}~--- вход: носитель, $m$ из~\eqref{eq:ucard}; выход: универсум мощности $m$: строится мультимножество \code{extra} с кратностями $\Kmax - 1$, из него выбирается \code{random\_submultiset} мощности $m - 2^n$, и к каждой кратности прибавляется 1; иначе \code{UniverseCardinalityError}.
\end{dashlist}

\subsection{Модуль ввода}
\label{sec:input}

Весь ввод читается построчно функцией \code{std::getline}~[6] и проверяется целиком. Оператор \code{>>} не подходит: он принимает «12abc» как 12, превращает «-1» в большое беззнаковое число и после ошибки оставляет поток в состоянии сбоя, из-за чего программа зацикливается.

\begin{dashlist}
\item \code{static std::string read\_line(prompt)}~--- печатает приглашение, читает строку, убирает пробелы, табуляции и \code{\textbackslash r} по краям; при конце потока~--- \code{InputClosedError}. Если ввод перенаправлен из файла (стандартный ввод не консоль), функция выводит прочитанную строку, чтобы протокол выглядел как сеанс в консоли.
\item \code{static bool parse\_number(text, value)}~--- истина, если строка состоит из 1--9 цифр; число записывается в \code{value}. Отвергаются знаки, точки, пробелы внутри и пустая строка.
\item \code{static bool is\_code(text, depth)}~--- истина, если строка длины $n$ из символов 0 и 1.
\item \code{size\_t read\_number(prompt, min, max)}~--- повторяет запрос с сообщением «Ошибка: введите целое число от min до max», пока не получит число из $[\text{min}, \text{max}]$.
\item \code{bool read\_number\_or\_empty(prompt, min, max, value)}~--- то же, но на пустую строку возвращает \code{false}.
\item \code{bool read\_code\_or\_empty(prompt, depth, code)}~--- повторяет запрос с сообщением «Ошибка: код должен состоять из $n$ символов 0 и 1», пока не получит код; на пустую строку возвращает \code{false}.
\end{dashlist}

\subsection{Класс Menu и функция main}

Класс \code{Menu} хранит состояние программы: общий носитель \code{carrier}, мультимножества \code{universe}, \code{a}, \code{b} и признаки \code{universe\_ready}, \code{a\_ready}, \code{b\_ready}. Признак создания универсума нужен отдельно от его пустоты: при $n = 0$ универсум создан, но пуст.

Меню выводит состояние ($n$, число элементов, $|U|$, $|A|$, $|B|$ или «не заполнено») и пункты: 1~--- создать $U$; 2, 3~--- заполнить $A$, $B$; 4~--- таблица кратностей $U$, $A$, $B$; 5~--- операции; 6~--- таблица кратностей всех тринадцати результатов; 0~--- выход. Пункты 2--4 требуют созданного универсума, пункты 5 и 6~--- ещё и заполненных $A$ и $B$; иначе выводится подсказка «Сначала создайте универсум (пункт 1).» или «Сначала заполните A и B (пункты 2 и 3).».

\textbf{Создание универсума} (\code{create\_universe}): читается $n \in [0, 19]$ и строится новый носитель. При $n = 0$ выводится «При n = 0 код Грея не содержит ни одного слова: U = \{\}, |U| = 0.», и способ задания кратностей не запрашивается, так как задавать нечего. Иначе читается способ (0~--- отмена, 1--3) и его параметры. Новые объекты собираются в локальных переменных и присваиваются полям только после успешного ввода, поэтому отмена на любом шаге сохраняет прежнее состояние. После создания $A$ и $B$ становятся пустыми мультимножествами над новым носителем со снятыми признаками; если они были заполнены, выводится сообщение, что они построены над прежним универсумом и очищены.

\textbf{Заполнение} (\code{fill\_multiset}): читается способ (0~--- отмена, 1~--- вручную, 2~--- автоматически), затем мощность $m \in [0, |U|]$; на пустую строку $m$ выбирается равновероятно из $[0, |U|]$ и выводится. Заполняется новый объект; поле и признак обновляются только при успехе. Автоматическое заполнение~--- \code{random\_submultiset(universe, m)}. Ручное заполнение (\code{fill\_manual}):
\begin{steps}
\item Вывести таблицу универсума.
\item Пока $|A| < m$: вывести остаток $d = m - |A|$; прочитать код $x$ (пустая строка~--- отмена всего заполнения).
\item Найти номер $x$ методом \code{find} и вычислить свободное место $f = \kU - \kA$. Если $f = 0$, вывести «Ошибка: элемент $x$ уже входит в A с наибольшей кратностью k\_U($x$) = …» и вернуться к шагу~2.
\item Прочитать добавляемую кратность $c \in [1, \min\{f, d\}]$, увеличить $\kA$ на $c$ и вернуться к шагу~2.
\end{steps}
Ограничение $c \leqslant \min\{f, d\}$ сохраняет $A \subseteq U$ и не даёт превысить заданную мощность; так как $m \leqslant |U|$, свободное место есть всегда, пока $|A| < m$. При $m = 0$ (в частности, при $n = 0$) цикл не выполняется ни разу и $A = \{\}$.

\textbf{Вывод.} \code{show\_table} печатает столбцы «№», «Код», «k\_U» (и «k\_A», «k\_B» с прочерком для незаполненного) и строку мощностей. \code{compute\_operations} вычисляет тринадцать результатов в виде списка «название, обозначение, мультимножество»; по нему \code{show\_operations} печатает строки «название обозначение = мультимножество, мощность $m$», а \code{show\_operations\_table}~--- таблицу кратностей всех результатов. Ширины столбцов считаются в символах, а не байтах: строки в UTF-8, и такие символы, как «∪», занимают несколько байт.

\textbf{Функция main} в Windows вызывает \code{SetConsoleOutputCP(CP\_UTF8)}, так как консоль по умолчанию выводит в кодировке 866; задаёт зерно (\code{--seed N} или \code{make\_seed}) и печатает его; запускает меню. \code{InputClosedError} обрабатывается как штатное завершение: «Ввод завершён.», код возврата 0. Любое другое исключение (возможно только при ошибке в программе)~--- сообщение и код возврата 1.

% ======================= 3 =======================
\section{Результаты работы}

Работа проверена на семи сценариях. Ввод каждого сценария записан в файл \code{report/scenarios/*.txt}, и программа запускалась командой \code{lab1 --seed N < файл}, где $N$~--- номер сценария; при повторном запуске вывод совпадает до символа. Так как ввод берётся из файла, введённые строки печатает сама программа (подраздел~\ref{sec:input}), и протокол совпадает с тем, что видно в консоли при ручном вводе. Рисунки получены из этих протоколов скриптом \code{report/make\_figures.py}: каждый рисунок начинается со строки выбора пункта меню; если рисунок охватывает несколько действий, меню между ними сохранено.

Кроме сценариев, ядро проверено отдельной программой: для $n = 0, \ldots, 10$ и 200 случайных пар $A$, $B$ на каждое $n$ все тринадцать результатов совпали с формулами таблицы~\ref{tab:ops} поэлементно и оказались подмультимножествами $U$; для $n = 1, \ldots, 16$ код Грея проверен на цикличность, различие соседних слов в одном разряде и совпадение с $i \oplus \lfloor i/2 \rfloor$; доля выбора элемента $x$ при $U = \{x^2, y^1\}$, $m = 1$ составила 0,6667 на $300\,000$ испытаний при теоретической $2/3$.

\subsection{Сценарий 1: случайный универсум, автоматическое заполнение}

После запуска программа печатает зерно и меню, в котором универсум не создан, а $A$ и $B$ не заполнены (рисунок~\ref{fig:fig_s1_start}).

\consolefig[0.6]{fig_s1_start}{Запуск программы, зерно 1}

Создаётся универсум с $n = 3$, способ 1, $K = 5$ (рисунок~\ref{fig:fig_s1_universe}). Коды идут в порядке таблицы~\ref{tab:gray}, кратности случайны в $[1, 5]$, $|U| = 26$.

\consolefig[0.6]{fig_s1_universe}{Создание универсума: $n = 3$, $K = 5$}

$A$ заполняется автоматически с заданной мощностью 10; для $B$ мощность не вводится (Enter) и выбирается случайно: $|B| = 12$ (рисунок~\ref{fig:fig_s1_fill}). Таблица пункта 4 (рисунок~\ref{fig:fig_s1_table}) подтверждает $\kA \leqslant \kU$ и $\kB \leqslant \kU$ для всех восьми элементов.

\consolefig[0.75]{fig_s1_fill}{Автоматическое заполнение: $|A|$ задана, $|B|$ выбрана случайно}

\consolefig[0.5]{fig_s1_table}{Таблица кратностей $U$, $A$, $B$}

Результаты операций~--- на рисунке~\ref{fig:fig_s1_ops}, таблица кратностей~--- на рисунке~\ref{fig:fig_s1_optable}.

\consolefig[1]{fig_s1_ops}{Результаты тринадцати операций (сценарий 1)}

\consolefig[1]{fig_s1_optable}{Таблица кратностей результатов (сценарий 1)}

Ручной расчёт для трёх элементов по таблице~\ref{tab:ops} приведён в таблице~\ref{tab:check1} и совпадает со строками 000, 010 и 100 рисунка~\ref{fig:fig_s1_optable}.

\begin{table}[H]
\caption{Ручная проверка (сценарий 1)}
\label{tab:check1}
\centering\small
\begin{tabularx}{\textwidth}{|l|C|C|C|}
\hline
 & 000: $k_U{=}3$, $k_A{=}3$, $k_B{=}2$ & 010: $k_U{=}5$, $k_A{=}3$, $k_B{=}1$ & 100: $k_U{=}5$, $k_A{=}1$, $k_B{=}4$ \\ \hline
$A \cup B$, $A \cap B$ & 3, 2 & 3, 1 & 4, 1 \\ \hline
$A \setminus B$ & $\min\{3, 1\} = 1$ & $\min\{3, 4\} = 3$ & $\min\{1, 1\} = 1$ \\ \hline
$B \setminus A$ & $\min\{2, 0\} = 0$ & $\min\{1, 2\} = 1$ & $\min\{4, 4\} = 4$ \\ \hline
$A \mathbin{\triangle} B$ & $\max\{1, 0\} = 1$ & $\max\{3, 1\} = 3$ & $\max\{1, 4\} = 4$ \\ \hline
$\overline{A}$, $\overline{B}$ & 0, 1 & 2, 4 & 4, 1 \\ \hline
$A + B$ & $\min\{5, 3\} = 3$ & $\min\{4, 5\} = 4$ & $\min\{5, 5\} = 5$ \\ \hline
$A - B$, $B - A$ & 1, 0 & 2, 0 & 0, 3 \\ \hline
$A \ast B$ & $\min\{6, 3\} = 3$ & $\min\{3, 5\} = 3$ & $\min\{4, 5\} = 4$ \\ \hline
$A / B$, $B / A$ & $\lfloor 3/2 \rfloor = 1$, $\lfloor 2/3 \rfloor = 0$ & 3, 0 & 0, 4 \\ \hline
\end{tabularx}
\end{table}

Элемент 000 показывает ограничение универсумом: без него $k_{A+B} = 5$ и $k_{A \ast B} = 6$, программа выводит 3. Элемент 010 показывает различие разностей: $k_{A \setminus B} = 3$, $k_{A - B} = 2$. Соотношения~\eqref{eq:check}: $|\overline{A}| = 26 - 10 = 16$, $|\overline{B}| = 26 - 12 = 14$, $|A \cup B| + |A \cap B| = 16 + 6 = 22 = 10 + 12$.

\subsection{Сценарий 2: ручной ввод с ошибками и отмена}

Универсум с $n = 2$ задаётся вручную (способ 2, рисунок~\ref{fig:fig_s2_universe}). Для элемента 00 вводятся 0, 101 и «abc»~--- все отвергаются с указанием диапазона $[1, 100]$; затем вводятся кратности 2, 1, 3, 1, и $U = \{00^2, 01^1, 11^3, 10^1\}$, $|U| = 7$.

\consolefig[0.65]{fig_s2_universe}{Ручное задание кратностей универсума с ошибками ввода}

$A$ заполняется вручную с мощностью 4 (рисунок~\ref{fig:fig_s2_fill_a}). Код «0» неверной длины и код «0a» с недопустимым символом отвергаются. Для 00 свободно $f = 2$, остаток $d = 4$, поэтому запрашивается кратность из $[1, 2]$; ввод 3 отвергается. Повторный ввод 00 отвергается: элемент уже взят с кратностью $k_U(00) = 2$. Для 11 свободно 3, но остаток 2, поэтому диапазон $[1, 2]$. Итог: $A = \{00^2, 11^2\}$. После каждой ошибки остаток не меняется, и ввод продолжается с того же места.

\consolefig[0.8]{fig_s2_fill_a}{Ручное заполнение $A$ с ошибками ввода}

Для $B$ выбирается ручное заполнение со случайной мощностью ($|B| = 3$), но на запрос кода вводится пустая строка: заполнение отменяется, и в меню $B$ по-прежнему «не заполнено». Затем $B$ заполняется автоматически с мощностью 3: $B = \{00^2, 11^1\}$ (рисунок~\ref{fig:fig_s2_fill_b}).

\consolefig[0.8]{fig_s2_fill_b}{Отмена ручного заполнения $B$ и автоматическое заполнение}

Результаты~--- на рисунках~\ref{fig:fig_s2_ops} и~\ref{fig:fig_s2_optable}. Для 00 ($k_U = k_A = k_B = 2$): $k_{A \setminus B} = \min\{2, 0\} = 0$, хотя элемент есть в $A$,~--- он целиком «занят» в $B$. Для 11 ($k_U = 3$, $k_A = 2$, $k_B = 1$): $k_{A \setminus B} = \min\{2, 2\} = 2$, $k_{A - B} = 1$, $k_{A+B} = \min\{3, 3\} = 3$, $k_{A/B} = 2$, $k_{B/A} = 0$.

\consolefig[0.85]{fig_s2_ops}{Результаты операций (сценарий 2)}

\consolefig[1]{fig_s2_optable}{Таблица кратностей результатов (сценарий 2)}

\subsection{Сценарий 3: универсум по мощности, граничные мощности $A$ и $B$}

Универсум с $n = 3$ задаётся по мощности (способ 3, рисунок~\ref{fig:fig_s3_universe}). Значения 7 и 801 вне отрезка $[2^3, 100 \cdot 2^3] = [8, 800]$ отвергаются, принимается $|U| = 20$; кратности распределены случайно, каждая не меньше 1, сумма равна 20.

\consolefig[0.6]{fig_s3_universe}{Универсум заданной мощности}

Мощность $|A| = 21 > |U|$ отвергается; при $|A| = 20 = |U|$ получается $A = U$. Для $B$ задана мощность 0: $B = \{\}$ (рисунок~\ref{fig:fig_s3_fill}).

\consolefig[0.8]{fig_s3_fill}{Граничные мощности: $A = U$, $B = \varnothing$}

Результаты (рисунки~\ref{fig:fig_s3_ops}, \ref{fig:fig_s3_optable}) совпадают с ожидаемыми для $A = U$, $B = \varnothing$: $A \cup B = A \setminus B = A \mathbin{\triangle} B = \overline{B} = A + B = A - B = U$; $A \cap B = B \setminus A = \overline{A} = B - A = A \ast B = \varnothing$. Оба частных пусты: $A / B$~--- так как все $k_B = 0$, $B / A$~--- так как $\lfloor 0 / k \rfloor = 0$.

\consolefig[1]{fig_s3_ops}{Результаты операций при $A = U$, $B = \varnothing$}

\consolefig[1]{fig_s3_optable}{Таблица кратностей при $A = U$, $B = \varnothing$}

\subsection{Сценарий 4: некорректный ввод в меню и при создании универсума}

До создания универсума выбираются пункты 2 и 5; программа выводит подсказку и возвращается в меню (рисунок~\ref{fig:fig_s4_order}).

\consolefig[0.55]{fig_s4_order}{Пункты 2 и 5 до создания универсума}

Как номер пункта вводятся «abc», 7, «-1», пустая строка, «1.5» и «2 3»; каждая строка отвергается сообщением «Ошибка: введите целое число от 0 до 6» (рисунок~\ref{fig:fig_s4_menu}).

\consolefig[0.5]{fig_s4_menu}{Некорректный ввод номера пункта}

Как разрядность вводятся 20 (больше $n_{\max} = 19$), «-2» и «abc»; как способ~--- 4; как $K$~--- 0 и 101. Всё отвергается с указанием допустимого диапазона (рисунок~\ref{fig:fig_s4_universe}). При $n = 3$, $K = 1$ все кратности равны 1: универсум~--- обычное множество, $|U| = 8$.

\consolefig[0.6]{fig_s4_universe}{Некорректный ввод разрядности, способа и $K$; универсум при $K = 1$}

Пункт 5 при незаполненных $A$ и $B$ выводит подсказку «Сначала заполните A и B (пункты 2 и 3).»; при заполнении $A$ способ 3 отвергается, а 0 отменяет заполнение (рисунок~\ref{fig:fig_s4_rest}).

\consolefig[0.55]{fig_s4_rest}{Пункт 5 до заполнения и отмена заполнения $A$}

\subsection{Сценарий 5: разрядность $n = 0$}
\label{sec:s5}

При $n = 0$ программа сообщает, что код Грея не содержит ни одного слова, выводит $U = \{\}$, $|U| = 0$ и не запрашивает способ задания кратностей; таблица не содержит строк (рисунок~\ref{fig:fig_s5_universe}). В меню универсум отмечен как созданный: «n = 0, элементов 0, |U| = 0 (U = \{\})».

\consolefig[0.7]{fig_s5_universe}{Создание универсума при $n = 0$}

$A$ заполняется вручную: мощность 1 отвергается (допустимый отрезок $[0, 0]$), при мощности 0 таблица пуста, цикл ручного ввода не выполняется ни разу, $A = \{\}$. $B$ заполняется автоматически со случайной мощностью; единственное возможное значение~--- 0, $B = \{\}$ (рисунок~\ref{fig:fig_s5_fill}).

\consolefig[0.8]{fig_s5_fill}{Ручное и автоматическое заполнение при $n = 0$}

Таблица пункта 4 содержит только строку мощностей $|U| = |A| = |B| = 0$ (рисунок~\ref{fig:fig_s5_table}). Все тринадцать результатов равны \code{\{\}} с мощностью 0, таблица результатов содержит только нулевую строку мощностей (рисунок~\ref{fig:fig_s5_ops}), что совпадает с подразделом~\ref{sec:universe}.

\consolefig[0.45]{fig_s5_table}{Таблица кратностей при $n = 0$}

\consolefig[0.95]{fig_s5_ops}{Результаты всех операций и таблица результатов при $n = 0$}

\subsection{Сценарий 6: пересоздание универсума и конец ввода}

Над универсумом $n = 2$, $K = 3$ ($|U| = 8$) заполняются $A = \{00^2, 01^2\}$ и $B = \{00^1, 01^1, 10^1\}$ (рисунок~\ref{fig:fig_s6_before}). Затем создаётся универсум $n = 3$, $K = 2$: программа сообщает, что $A$ и $B$ были построены над прежним универсумом и очищены, и выводит новую таблицу из восьми элементов, $|U| = 12$ (рисунок~\ref{fig:fig_s6_recreate}). Мультимножества старого носителя не могут попасть в операции с новым универсумом: их признаки сняты, а при попытке смешать носители ядро выбросило бы \code{CarrierMismatchError}.

\consolefig[0.6]{fig_s6_before}{Заполнение $A$ и $B$ над универсумом с $n = 2$}

\consolefig[0.6]{fig_s6_recreate}{Пересоздание универсума с $n = 3$}

Начинается ручное заполнение $A$ мощности 2, и на запросе кода поток ввода заканчивается (в консоли~--- Ctrl+Z и Enter, при вводе из файла~--- конец файла). Программа выводит «Ввод завершён.» и завершается с кодом 0, не зацикливаясь на повторных запросах (рисунок~\ref{fig:fig_s6_eof}).

\consolefig[0.8]{fig_s6_eof}{Конец ввода во время ручного заполнения}

\subsection{Сценарий 7: наибольшая разрядность $n = 19$}

Создаётся универсум $n = 19$, $K = 100$: $2^{19} = 524\,288$ элементов, $|U| = 26\,489\,557$. Таблица обрезана после 32 строк с пометкой (рисунок~\ref{fig:fig_s7_universe}). $A$ заполняется автоматически со случайной мощностью $|A| = 21\,290\,598$; запись $A$ обрезана после 16 элементов с указанием числа скрытых (рисунок~\ref{fig:fig_s7_fill}). Весь сценарий, включая заполнение и вывод, выполняется за доли секунды, что согласуется с таблицей~\ref{tab:bench}.

\consolefig[0.6]{fig_s7_universe}{Универсум при $n = 19$, $K = 100$}

\consolefig[1]{fig_s7_fill}{Автоматическое заполнение $A$ при $n = 19$}

% ======================= ЗАКЛЮЧЕНИЕ =======================
\structsection{ЗАКЛЮЧЕНИЕ}

Разработана консольная программа на C++, которая строит бинарный код Грея разрядности от 0 до 19 по номеру инвертируемого разряда, задаёт на нём универсум-мультимножество одним из трёх способов (случайные кратности от 1 до $K$, ручной ввод кратностей, заданная или случайная мощность), заполняет $A$ и $B$ вручную или автоматически с заданной или случайной мощностью от 0 до $|U|$ и выводит тринадцать результатов операций, таблицы кратностей и мощности. Работа проверена на семи воспроизводимых сценариях: результаты совпали с ручным расчётом (таблица~\ref{tab:check1}), с ожидаемыми значениями в граничных случаях $A = U$, $B = \varnothing$, $n = 0$, $K = 1$, а все некорректные значения отвергнуты с сообщением и повторным запросом.

\textbf{Достоинства и недостатки.}
\begin{dashlist}
\item Генерация кода Грея не использует рекурсию и не хранит код меньшего порядка: каждое слово получается из предыдущего инвертированием одного разряда, всего $O(2^n)$ действий. Недостаток: слова выдаются только подряд; получить слово по номеру без предыдущих шагов алгоритм не позволяет (это делает формула $i \oplus \lfloor i/2 \rfloor$).
\item Мультимножество хранится функцией кратности по номерам носителя: операции выполняются за $O(2^n)$ без поиска, универсум~--- объект того же класса, а результат любой операции остаётся подмультимножеством $U$ (утверждение подраздела~1.3). Цена~--- память $O(2^n)$ на каждое мультимножество, даже почти пустое.
\item Поиск элемента по коду выполняется за $O(n)$ через обратную таблицу \code{positions} ценой ещё $2^n$ чисел памяти.
\item Случайное заполнение даёт точно заданную мощность, все наборы «единиц кратности» равновероятны, холостых выборок нет. Недостаток: время $O(|U|)$, а не $O(2^n)$; именно оно ограничивает $n_{\max} = 19$.
\item Ручной ввод универсума требует $2^n$ чисел и практичен только при малых $n$; при больших $n$ вывод обрезается, и весь результат на экране не увидеть.
\item При $n = 0$ принято $U = \{\}$, тогда как формально множество слов длины 0 состоит из одного пустого слова; соглашение выбрано потому, что пустое слово нельзя ни показать, ни ввести.
\end{dashlist}

\textbf{Масштабирование.} Функциональность расширяется незначительными изменениями кода, так как математическая часть не зависит от интерфейса, все операции выражены через кратности по номерам, а меню строит вывод и таблицы по общему списку операций.
\begin{dashlist}
\item \textit{Новая операция} (умножение мультимножества на число $\lambda$: $\min\{\lambda \kA, \kU\}$; проверка включения $A \subseteq B$ или равенства $A = B$): один метод \code{Multiset} с одним циклом по кратностям (для проверок уже есть \code{is\_submultiset\_of}) и одна строка в \code{Menu::compute\_operations}; пункты 5 и 6 подхватят её автоматически.
\item \textit{Больше двух мультимножеств}: поля \code{a}, \code{b} заменяются массивом мультимножеств с именами и признаками; \code{fill\_multiset} уже получает мультимножество, признак и имя параметрами, поэтому меняются только пункты меню и цикл в \code{compute\_operations}.
\item \textit{Другой носитель} (обычный двоичный код, код Грея с другим порядком, произвольные слова): меняется только конструктор \code{CodeSet}; операции работают с номерами элементов и носитель не используют.
\item \textit{Вывод в файл}: \code{print\_multiset} уже пишет в произвольный \code{std::ostream}; достаточно передать поток в \code{show\_table} и \code{show\_operations} вместо \code{std::cout}.
\item \textit{Пакетный режим}: уже работает~--- \code{lab1 --seed N < файл} выполняет сценарий и печатает протокол; так получены все рисунки раздела~3.
\item \textit{Большие $n$}: узкое место~--- \code{random\_submultiset}. Если вместо просмотра каждой единицы выбирать $\kA$ сразу из гипергеометрического распределения, время станет $O(2^n)$; по замеру \code{tools/bench.cpp} операции и построение универсума способом~1 укладываются в 1~с при всех измеренных $n$ до 22 включительно (при $n = 22$~--- 264 и 47~мс). Меняются функция \code{random\_submultiset} и константа \code{MAX\_DEPTH}; при $n > 31$ нужна 64-разрядная сборка (\code{MAX\_TYPE\_DEPTH} вычисляется из размера \code{size\_t} автоматически).
\item \textit{Большие кратности}: константа \code{MAX\_MULTIPLICITY}; при этом нужно повторить замер (время автозаполнения пропорционально $\Kmax$) и сохранить $\Kmax^2 < 2^{32}$, чтобы не переполнялось произведение кратностей в 32-разрядной сборке.
\end{dashlist}

В ходе работы освоены построение кода Грея по номеру изменяемого разряда с доказательством корректности, представление мультимножества функцией кратности и вычисление операций над ним, равновероятный выбор подмультимножества заданной мощности алгоритмом выборки, обоснование ограничений программы замером, защищённый построчный ввод и разделение программы на математическую часть и интерфейс.

% ======================= ИСТОЧНИКИ =======================
\structsection{СПИСОК ИСПОЛЬЗОВАННЫХ ИСТОЧНИКОВ}

\begin{enumerate}[label=\arabic*, leftmargin=*, labelindent=\parindent, itemsep=0pt, align=left]
\item Новиков, Ф.\,А. Дискретная математика для программистов : учебник для вузов / Ф.\,А.~Новиков.~--- 3-е изд.~--- Санкт-Петербург : Питер, 2009.~--- 384~с.
\item Кнут, Д.\,Э. Искусство программирования. Т.~2. Получисленные алгоритмы / Д.\,Э.~Кнут.~--- 3-е изд.~--- Москва : Вильямс, 2000.~--- 832~с.
\item Nielsen, J. Usability Engineering / J.~Nielsen.~--- Boston : Academic Press, 1993.~--- 362~p.
\item Востров, А.\,В. Дискретная математика : курс лекций / А.\,В.~Востров ; Санкт-Петербургский политехнический университет Петра Великого.~--- Санкт-Петербург, 2026.~--- URL: \url{https://tema.spbstu.ru/dismath_lect/} (дата обращения: 04.10.2026).~--- Текст : электронный.
\item std::uniform\_int\_distribution // cppreference.com : [сайт].~--- URL: \url{https://en.cppreference.com/w/cpp/numeric/random/uniform_int_distribution} (дата обращения: 04.10.2026).~--- Текст : электронный.
\item std::getline // cppreference.com : [сайт].~--- URL: \url{https://en.cppreference.com/w/cpp/string/basic_string/getline} (дата обращения: 04.10.2026).~--- Текст : электронный.
\end{enumerate}

% ======================= ПРИЛОЖЕНИЕ =======================
\clearpage
\phantomsection
\addcontentsline{toc}{section}{ПРИЛОЖЕНИЕ А Текст программы}
\begin{center}
\textbf{ПРИЛОЖЕНИЕ А}\\
\textbf{Текст программы}
\end{center}

\codefile{\code{CMakeLists.txt}}{lst:CMakeListstxt}{../CMakeLists.txt}
\codefile{\code{src/core/Random.hpp}}{lst:srccoreRandomhpp}{../src/core/Random.hpp}
\codefile{\code{src/core/Random.cpp}}{lst:srccoreRandomcpp}{../src/core/Random.cpp}
\codefile{\code{src/core/GrayCode.hpp}}{lst:srccoreGrayCodehpp}{../src/core/GrayCode.hpp}
\codefile{\code{src/core/GrayCode.cpp}}{lst:srccoreGrayCodecpp}{../src/core/GrayCode.cpp}
\codefile{\code{src/core/Multiset.hpp}}{lst:srccoreMultisethpp}{../src/core/Multiset.hpp}
\codefile{\code{src/core/Multiset.cpp}}{lst:srccoreMultisetcpp}{../src/core/Multiset.cpp}
\codefile{\code{src/core/Universe.hpp}}{lst:srccoreUniversehpp}{../src/core/Universe.hpp}
\codefile{\code{src/core/Universe.cpp}}{lst:srccoreUniversecpp}{../src/core/Universe.cpp}
\codefile{\code{src/ui/Input.hpp}}{lst:srcuiInputhpp}{../src/ui/Input.hpp}
\codefile{\code{src/ui/Input.cpp}}{lst:srcuiInputcpp}{../src/ui/Input.cpp}
\codefile{\code{src/ui/Menu.hpp}}{lst:srcuiMenuhpp}{../src/ui/Menu.hpp}
\codefile{\code{src/ui/Menu.cpp}}{lst:srcuiMenucpp}{../src/ui/Menu.cpp}
\codefile{\code{src/main.cpp}}{lst:srcmaincpp}{../src/main.cpp}
\codefile{\code{tools/bench.cpp}}{lst:toolsbenchcpp}{../tools/bench.cpp}

\end{document}
